\documentclass[10pt,a4paper]{amsart}
\usepackage[margin=19mm]{geometry}
\usepackage{amsmath,amssymb,mathtools}
\usepackage[scr=rsfs,cal=euler]{mathalfa}
\usepackage{xfrac}
\usepackage[unicode,hidelinks,bookmarks=false]{hyperref}
\newcommand{\ie}{i.e.\ }
\newcommand{\cf}{cf.\ }
\newcommand{\resp}{respectively\ }
\newcommand{\Subsection}{\subsection}
\providecommand{\nobreakdash}{\nobreak\hbox{-}}
\setlength{\emergencystretch}{2em}
\multlinegap0pt
\usepackage{bm}
\usepackage{bbm}
\usepackage{dsfont}
\usepackage{xspace}
\usepackage{cancel}
\usepackage{mathtools}
\usepackage{amsthm}
\makeatletter
\@ifundefined{lemma}{\newtheorem{lemma}{Lemma}}{}
\@ifundefined{theorem}{\newtheorem{theorem}{Theorem}}{}
\makeatother
\def\R{\mathbb{R}}
\def\vu{ u }
\title[Liouville Theorems for Stationary Navier-Stokes Equations vi]{Liouville Theorems for Stationary Navier-Stokes Equations via the Radial Velocity Component}
\author{Gaston Vergara-Hermosilla}
\date{Source: arXiv:2605.05647v1 (CC BY 4.0)}
\begin{document}
\maketitle

\noindent\emph{Source and licence.} Liouville Theorems for Stationary Navier-Stokes Equations via the Radial Velocity Component. Version arXiv:2605.05647v1 (CC BY 4.0), distributed under the Creative Commons Attribution 4.0 International licence, as recorded in the arXiv licence metadata for this exact version and shown on the abstract page https://arxiv.org/abs/2605.05647v1. Full rights notice: licenses/arxiv-2605.05647-CC-BY-4.0.txt.\par\smallskip

\noindent\textit{Editorial scope.} Counted unit: the Lemma whose source label is \texttt{prop:main\_limit\_thmover\_new} in the article (source0.tex lines 482--492, source label \texttt{prop:main\_limit\_thmover\_new}), delivered together with the article's own complete original proof of it (source0.tex lines 493--597, one \texttt{proof} environment of 105 lines). No mathematics is written, repaired or re-derived: statement and proof are the article's own text line for line. Required context retained uncounted: SNS (lines 210--218); thm1.2 (lines 301--308). Every definition, display and label of those lines is retained unchanged. Omitted from this package and not redistributed: 1--209, 219--300, 309--481, 598--1124. Nothing inside a retained excerpt is omitted or altered: every retained body is delivered exactly as the article's lines are written, so no omission receipt of any kind accompanies this package. Citations: the retained excerpts carry no citation command at all, so no bibliography entry of the article is transcribed and no reference is redistributed. Reference adaptation: none was needed and none was made; every reference inside the delivered unit resolves inside it, so no unresolved reference or citation is produced. Printed slips kept literal: no printed word, symbol, hypothesis or number of the counted statement or of its proof was altered. Layout and class: the article's own class is \texttt{amsart} with packages including setspace, geometry, amsmath, amsthm, amsfonts, amssymb, mathtools, mathrsfs, url, bm, enumitem, stackengine, newtxtext, newtxmath; the wrapper is the lane's pinned \texttt{amsart} editorial wrapper at 10pt on A4, which restates the theorem-like environments the retained text needs and adds the article's own macro definitions that the retained text needs (\texttt{\textbackslash R}, \texttt{\textbackslash vu}), with extra wrapper packages (bm, bbm, dsfont, xspace, cancel, mathtools). The delivered numbering is the unit's own. Assets: no retained span contains a figure environment, a graphics-inclusion command, a PDF-page inclusion, a table or a TikZ picture; no image file is copied into this package, the package declares no asset dependency and no image file is redistributed with it. Licence: version arXiv:2605.05647v1 of this article is distributed under the Creative Commons Attribution 4.0 International licence, as recorded in the arXiv licence metadata for this exact version and shown on the abstract page https://arxiv.org/abs/2605.05647v1. A full rights notice is delivered at licenses/arxiv-2605.05647-CC-BY-4.0.txt. Characters: every retained excerpt is pure ASCII, so the delivered engine, XeLaTeX with its default Latin Modern text font, reports no missing character.
	\begin{equation}\label{SNS}
		\begin{cases}
			-\Delta u+  ( u\cdot \nabla ) u +\nabla \pi  =0,
			\\
			\quad \nabla \cdot u=0, 
		\end{cases}
		\quad 
		\text{in }\R^3,
	\end{equation} 

  \begin{theorem}\label{thm1.2} 
 Let $\vu  \in \dot{H}^1(\mathbb{R}^3) $ be a solution of (\ref{SNS})
 and let $u_\rho(x)$ be the radial component of such solution in spherical coordinates. 
Consider $R_0>3$ fixed  and let $\eta(\cdot)$ be a scalar function defined by  $\eta(x)=  3 $ for $|x|< R_0$ and $\eta(x)=   \frac{3R_0}{|x|},$ for $|x| \geq  R_0$.  
 Suppose in addition that
 $u_\rho \in L^{3 + \eta(\cdot) }(\mathbb{R}^3)$.
 Then, $u $ is identically equal to zero.
\end{theorem} 

\begin{lemma}\label{prop:main_limit_thmover_new}
Consider $p(\cdot)$ be the variable exponent defined in Theorem \ref{thm1.2}, $R_0>3$ and $R\geq 2R_0^2$. Let  
$h:[1,\infty) \to [0,\infty)$ be such that 
$
h(R) \to 0  \text{ as } R \to \infty.
$ 
Then
\[
\lim_{R\to\infty} R^{1 - \frac{3}{p_{\mathcal{C}_R}^+}} h(R) = 0.
\]
\end{lemma}

\begin{proof} 
Let $x\in\mathcal{C}( R) $. 
Since $p(\cdot)$ is continuous and radially decreasing, there exists a decreasing and continuous function
$\tilde{p} : [0, \infty) \to [3, 6]$ such that $ p(x) = \tilde{p}(|x|) $ for all $ x \in \mathbb{R}^3 $.
For each $ x \in\mathcal{C}( R)  $, we have $|x| > R/2$, hence
\[
p(x) = \tilde{p}(|x|) \leq \tilde{p}(R/2),
\]
and therefore we can  write
\[
\operatorname{ess\,sup}_{x \in\mathcal{C}( R) } p(x) \leq \tilde{p}(R/2).
\]
Consider $\varepsilon > 0$. By continuity of $\tilde{p}$ at $R/2$, there exists $\delta > 0$ such that
\[
r \in (R/2, R/2 + \delta) \;\Rightarrow\; |\tilde{p}(r) - \tilde{p}(R/2)| < \varepsilon,
\]
which implies $\tilde{p}(r) > \tilde{p}(R/2) - \varepsilon$.
To continue, we define
\[
A(\delta,R) := \{x \in \mathbb{R}^3 : R/2 < |x| < R/2 + \delta\}.
\]
Then $A(\delta,R) \subset\mathcal{C}( R) $, and for all $x \in A_\delta$, we have
\[
p(x) > \tilde{p}(R/2) - \varepsilon.
\]
Moreover, since set $A(\delta,R)$ has positive measure,   we can write
\[
\operatorname{ess\,sup}_{x \in\mathcal{C}( R) } p(x) \geq \tilde{p}(R/2) - \varepsilon.
\]
Since $\varepsilon>0$ is arbitrary, we conclude
$
p_{\mathcal{C}_R}^+ = \tilde{p}(R/2).
$ 
Now since $R \geq 2R_0^2$, we have $R/2 \geq R_0^2$, and by assumption, we have 
\[
\left| \tilde p(R/2) - 3 \right| \leq \frac{C}{R/2} = \frac{2C}{R}<1.
\]
Defining $\varepsilon_R := \tilde p(R/2) - 3$, we conclude  
that there exists a sequence $(\varepsilon_R)$ such that
\[
p_{\mathcal{C}_R}^+ = 3 + \varepsilon_R, 
\quad 0<\varepsilon_R \le \frac{2C}{R}<1.
\]
In particular, note that 
\[
p_{\mathcal{C}_R}^+ = 3 + O\!\left( \frac{1}{R}\right).
\]

To continue, define $
\omega_R := 1 - \frac{3}{p_{\mathcal{C}_R}^+}$.   Thus, we can write 
\[
\frac{3}{p_{\mathcal{C}_R}^+}
= \frac{3}{3 + \varepsilon_R}
= \frac{1}{1 + \frac{\varepsilon_R}{3}}.
\]
Now, by using the Taylor expansion for $t = \varepsilon_R/3 < 1$, we get 
\[
\frac{1}{1+t} = 1 - t + O(t^2),
\]
and then we conclude 
\[
\frac{3}{p_{\mathcal{C}_R}^+}
= 1 - \frac{\varepsilon_R}{3} + O(\varepsilon_R^2).
\]
Therefore, we can write 
\[
\omega_R
= 1 - \frac{3}{p_{\mathcal{C}_R}^+}
= \frac{\varepsilon_R}{3} + O(\varepsilon_R^2)
= O\!\left(\frac{1}{R}\right).
\]
Now, using the identity 
$
R^{\omega_R} = \exp(\omega_R \ln R),
$ 
and the fact that $\omega_R = O(1/R)$, we get
\[
\omega_R \ln R = O\!\left(\frac{\ln R}{R}\right) \to 0.
\]
Thus, considering $s:=\omega_R \ln R$ and using the expansion $e^s = 1 + O(s)$ as $s \to 0$, we conclude
\[
R^{\omega_R} = 1 + O\!\left(\frac{\ln R}{R}\right).
\]
Provided with this information   we can write
\[
R^{1 - \frac{3}{p_{\mathcal{C}_R}^+}} = R^{\omega_R}
= 1 + \alpha_R,
\]
where $\alpha_R = O\!\left(\frac{\ln R}{R}\right)$.
Thus, we obtain 
\[
R^{1 - \frac{3}{p_{\mathcal{C}_R}^+}} h(R)
= (1 + \alpha_R) h(R)
= h(R) + \alpha_R h(R).
\]
Since $h(R) \to 0$ and $(\alpha_R)$ is bounded for $R \geq 1$, we have
\[
|\alpha_R h(R)| \le C |h(R)| \to 0 \quad \text{as } R \to \infty.
\]
Thus, conclude 
$ 
R^{1 - \frac{3}{p_{\mathcal{C}_R}^+}} g(R) \to 0   \text{ as } R \to \infty,
$ 
and we finish the proof. 
\end{proof}

\end{document}
